\section{FiLM Rate Conditioning Head}
\label{sec:film-head}

\subsection{Rate Representation}

The compression rate is represented as:
\begin{equation}
  r = \log_2(\bpp),
\end{equation}
where $\bpp$ is the bits per point of the compressed bitstream.
The logarithmic transformation linearises the rate-distortion relationship:
distortion typically scales as $2^{-\alpha r_{\text{linear}}}$ in bits per point,
so $\log_2(\bpp)$ provides a scale on which equal intervals correspond to
approximately equal changes in distortion.

For the training set (r2, r4, r6 rate points), typical $\bpp$ values are
approximately 0.04, 0.025, and 0.012, giving $r$ values of approximately
$-4.6$, $-5.3$, and $-6.4$.
The logarithmic scale spaces these more uniformly than the raw $\bpp$ values.

\subsection{Rate MLP}

A two-layer MLP maps the scalar $r$ to the FiLM parameters:
\begin{equation}
  \mathbf{e}(r) = \mathbf{W}_2\,\sigma(\mathbf{W}_1 r + \mathbf{b}_1) + \mathbf{b}_2
  \in \R^{64}, \qquad \sigma = \text{ReLU},
\end{equation}
with $\mathbf{W}_1 \in \R^{64 \times 1}$, $\mathbf{W}_2 \in \R^{64 \times 64}$,
$\mathbf{b}_1 \in \R^{64}$, $\mathbf{b}_2 \in \R^{64}$.
The 64-dimensional embedding is split into the scale and shift vectors:
$\boldsymbol{\gamma} = \mathbf{e}(r)_{1:32}$ and
$\boldsymbol{\beta} = \mathbf{e}(r)_{33:64}$, each of dimension 32.

The MLP adds $64 + 64 \times 64 + 64 = 4\,224$ parameters,
approximately 0.17\% of the total model size.

\subsection{Feature Modulation}

The modulation is applied to every active voxel's feature vector $\mathbf{h}_i \in \R^{32}$
at the first decoder level output:
\begin{equation}
  \mathbf{h}'_i = \boldsymbol{\gamma}(r) \odot \mathbf{h}_i + \boldsymbol{\beta}(r),
  \qquad \forall i \in \{1, \ldots, N\}.
\end{equation}
The scaling $\boldsymbol{\gamma}(r)$ applies a per-channel amplitude rescaling
that can amplify or suppress specific feature dimensions depending on the rate.
The shift $\boldsymbol{\beta}(r)$ adds a per-channel bias, effectively
adjusting the ``resting state'' of the features at each rate.

\paragraph{What FiLM learns.}
At r2, where corrections are large, the FiLM layer is expected to amplify
feature channels that detect surface discontinuities and suppress channels
related to confident local geometry.
At r6, where corrections are small, the amplitude should be lower overall,
and the pattern of active channels may differ.
The MLP can learn any continuous function of $r$ within its capacity,
allowing rate-specific feature representations without storing separate weights
per rate.

\subsection{Displacement Head}

The modulated 32-channel feature maps are passed through a final $1^3$ sparse
convolution (same-resolution, no submanifold constraint needed here since we
output per-active-coordinate values), followed by $\tanh$ scaling:
\begin{equation}
  \Delta_i = 5 \cdot \tanh\!\left(\mathbf{W}_{\text{head}} \mathbf{h}'_i + \mathbf{b}_{\text{head}}\right),
  \qquad \mathbf{W}_{\text{head}} \in \R^{3 \times 32},\; \mathbf{b}_{\text{head}} \in \R^3.
\end{equation}

The output displacement is added to the integer input coordinate to
produce the refined position in continuous $\R^3$:
\begin{equation}
  \hat{p}^{\text{ref}}_i = \hat{p}_i + \Delta_i.
\end{equation}
Note that the refined positions are in $\R^3$, not $\mathbb{Z}^3$:
the output point cloud is a floating-point cloud, not a voxelized integer grid.
For codec comparison and metric computation, the refined cloud is compared
directly with the original (which is also in integer voxel coordinates)
using the standard nearest-neighbour PSNR computation.

% ============================================================
